% 分布式 SGD 拓扑对比：参数服务器（左）vs Ring-AllReduce（右）
% 箭头表示梯度/参数通信方向
\documentclass[tikz,border=4pt]{standalone}
\usepackage{ctex}
\usepackage{amsmath}
\usetikzlibrary{calc, arrows.meta, positioning, shapes.geometric, fit}

\begin{document}
\begin{tikzpicture}[scale=1.0, thick,
  server/.style={draw=red!70!black, fill=red!15, rounded corners=4pt,
                 minimum width=1.6cm, minimum height=0.7cm, font=\footnotesize, align=center},
  worker/.style={draw=blue!70!black, fill=blue!12, rounded corners=4pt,
                 minimum width=1.4cm, minimum height=0.6cm, font=\footnotesize, align=center},
  myarrow/.style={-{Stealth[length=5pt]}, very thick},
  bidir/.style={{Stealth[length=5pt]}-{Stealth[length=5pt]}, very thick},
]

% ══════════════════════════════════════════════
% 左半部分：参数服务器架构
% ══════════════════════════════════════════════

% 标题
\node[font=\small\bfseries, red!70!black] at (2.0, 5.2) {参数服务器（Parameter Server）};

% 中央参数服务器
\node[server, minimum width=2.0cm] (PS) at (2.0, 3.8) {参数服务器\\PS};

% 四个 Worker 节点
\node[worker] (W1) at (0.0, 2.0) {Worker 1};
\node[worker] (W2) at (1.3, 2.0) {Worker 2};
\node[worker] (W3) at (2.7, 2.0) {Worker 3};
\node[worker] (W4) at (4.0, 2.0) {Worker 4};

% PS 到 Worker 的连线（下发参数：蓝色实线向下）
\foreach \w in {W1, W2, W3, W4}{
  \draw[myarrow, blue!70, bend left=8] (PS.south) to (\w.north);
}

% Worker 到 PS 的连线（上传梯度：橙色虚线向上）
\draw[myarrow, orange!80!black, dashed, bend right=8] (W1.north) to (PS.south);
\draw[myarrow, orange!80!black, dashed, bend right=8] (W2.north) to (PS.south);
\draw[myarrow, orange!80!black, dashed, bend right=8] (W3.north) to (PS.south);
\draw[myarrow, orange!80!black, dashed, bend right=8] (W4.north) to (PS.south);

% 图例（左下角）
\draw[myarrow, blue!70] (0.0, 0.9) -- (0.9, 0.9);
\node[right, font=\tiny, blue!70] at (0.9, 0.9) {下发参数};
\draw[myarrow, orange!80!black, dashed] (0.0, 0.5) -- (0.9, 0.5);
\node[right, font=\tiny, orange!80!black] at (0.9, 0.5) {上传梯度};

% PS 瓶颈标注
\node[font=\tiny, red!60, align=center] at (2.0, 3.2)
  {带宽瓶颈：$O(Nd)$};

% ══════════════════════════════════════════════
% 右半部分：Ring-AllReduce 环形拓扑
% ══════════════════════════════════════════════

% 标题
\node[font=\small\bfseries, blue!70!black] at (7.8, 5.2) {Ring-AllReduce（环形）};

% 5 个 Worker 节点，排成正五边形环
% 中心 (7.8, 3.2)，半径 1.5
\def\cx{7.8}
\def\cy{3.2}
\def\r{1.55}

% 节点位置（从顶部顺时针）
\node[worker] (R1) at (\cx,           {\cy+\r})       {GPU 0};
\node[worker] (R2) at ({\cx+\r*0.951}, {\cy+\r*0.309}) {GPU 1};
\node[worker] (R3) at ({\cx+\r*0.588}, {\cy-\r*0.809}) {GPU 2};
\node[worker] (R4) at ({\cx-\r*0.588}, {\cy-\r*0.809}) {GPU 3};
\node[worker] (R5) at ({\cx-\r*0.951}, {\cy+\r*0.309}) {GPU 4};

% 环形箭头（顺时针，蓝色）
\draw[bidir, blue!60] (R1.south east) -- (R2.north west);
\draw[bidir, blue!60] (R2.south)      -- (R3.north);
\draw[bidir, blue!60] (R3.west)       -- (R4.east);
\draw[bidir, blue!60] (R4.north)      -- (R5.south);
\draw[bidir, blue!60] (R5.north east) -- (R1.south west);

% 通信量标注
\node[font=\tiny, blue!70!black, align=center] at (\cx, \cy)
  {$\approx 2d$ 通信量\\（与 $N$ 无关）};

% 无瓶颈标注
\node[font=\tiny, green!50!black, align=center] at (\cx, 0.7)
  {均匀负载，无单点瓶颈};

% ══════════════════════════════════════════════
% 分割线
% ══════════════════════════════════════════════
\draw[gray!40, very thick, dashed] (5.5, 0.2) -- (5.5, 5.5);

% ══════════════════════════════════════════════
% 底部标题
% ══════════════════════════════════════════════
\node[font=\small\bfseries] at (5.8, -0.2)
  {分布式 SGD 通信拓扑对比};

\end{tikzpicture}
\end{document}
